% --- Archivo: exercises.tex ---

\section*{Ejercicios del Capítulo}
\addcontentsline{toc}{section}{Ejercicios del Capítulo}

\begin{exercise}\label{v2:ejer:1.5.1}
    Probar el Lema \ref{v2:lem:1.5.3} utilizando la técnica de inducción.
\end{exercise}

\begin{exercise}\label{v2:ejer:1.6.1}
    Probar que si $F : \Rn \to \R^l$ es Lipschitz-continua con módulo $L > 0$ en una vecindad del punto $x \in \Rn$, entonces $\partial F(x)$ es un conjunto compacto no vacío y $\|H\| \le L \ \forall H \in \partial F(x)$.
\end{exercise}

\begin{exercise}\label{v2:ejer:1.6.2}
    Probar que si la función $F : \Rn \to \R$ es convexa en una vecindad (convexa) del punto $x \in \Rn$, entonces $F$ es semi-diferenciable en el punto $x$.
\end{exercise}

\begin{exercise}\label{v2:ejer:1.6.3}
    Probar que si las funciones $F_1, F_2 : \Rn \to \R^l$ son semi-diferenciables en el punto $x \in \Rn$, entonces las funciones definidas a continuación también son semi-diferenciables en $x$:
    \begin{enumerate}[(a)]
        \item $F : \Rn \to \R$, $F(x) = \interno{F_1(x)}{F_2(x)}$.
        \item $F : \Rn \to \R^l$, $F(x) = F_1(x) + F_2(x)$.
        \item $F : \Rn \to \R^l$, $F_i(x) = \max\{(F_1(x))_i, (F_2(x))_i\}$, $i = 1, \dots, l$.
        \item $F : \Rn \to \R^l$, $F_i(x) = \min\{(F_1(x))_i, (F_2(x))_i\}$, $i = 1, \dots, l$.
    \end{enumerate}
\end{exercise}

\begin{exercise}\label{v2:ejer:1.6.4}
    Probar que si las funciones $F_1 : \Rn \to \R^{l_1}$ y $F_2 : \Rn \to \R^{l_2}$ son semi-diferenciables en el punto $x \in \Rn$, entonces la función $F : \Rn \to \R^{l_1} \times \R^{l_2}$, $F(x) = (F_1(x), F_2(x))$, es semi-diferenciable en $x$.
\end{exercise}

\begin{exercise}\label{v2:ejer:1.6.5}
    Probar que si la función $F_1 : \Rn \to \R^{l_1}$ es semi-diferenciable en el punto $x \in \Rn$, y la función $F_2 : \R^{l_1} \to \R^{l_2}$ es semi-diferenciable en el punto $F_1(x)$, entonces la función $F : \Rn \to \R^{l_2}$, $F(x) = F_2(F_1(x))$ es semi-diferenciable en $x$.
\end{exercise}
%%% Local Variables:
%%% mode: LaTeX
%%% TeX-master: "../../../Apuntes-Programacion-No-Lineal"
%%% End:
